%%%%%%%%%%
%% 2025 %%
%%%%%%%%%%

@article{AbbotMalani2025,
   author = {Abbot, Dorian S. and Malani, Anup},
   title = {Revisiting the physical limits to economic growth, with a focus on the waste heat limit},
   journal = {PLOS ONE},
   volume = {20},
   number = {3},
   pages = {e0319217},
   abstract = {Whether there are imminent physical limits to economic growth has been the focus of a long-standing debate, often (but not exclusively) pitting natural scientists against economists. One side of the debate posits that finite resources constrain human economic growth, the other that productivity improvements and substitution can relax those resource constraints. However, the two sides lack a common framework to resolve their disagreement. We provide this framework and mediate this dispute using a combination of low-order economic and physical models, focusing especially on the waste heat limit. To start, we use a simple accounting identity to show that historical rates of growth in average productivity can support modest economic and population growth (totaling 1.2\%). We model the future trajectory of productivity growth by combining a canonical production function (Cobb-Douglas) with a flexible model where population growth drives total-factor-productivity growth. We find remarkable parameter sensitivity: across a plausible range of parameters waste heat either never constrains economic growth in the model or does so in a few hundred years. Interestingly, either decreasing or increasing population growth can dramatically extend the time to heat constraint in the model, the latter due to increases in innovation. Finally, we use a more flexible (constant elasticity of substitution) production function and historical changes in the relative use of different resources to show that humans have the ability to substitute away from resources in shorter supply. Although there are vast uncertainties in this type of projection, our work highlights the potential for an optimistic future for humanity in which economic growth continues on a millennial timescale.},
   DOI = {10.1371/journal.pone.0319217},
   url = {https://doi.org/10.1371/journal.pone.0319217},
   year = {2025},
   type = {Journal Article}
}

% Working papers

@techreport{Leder-LuisMalani2025,
 title = "The Economics of Healthcare Fraud",
 author = "Leder-Luis, Jetson and Malani, Anup",
 institution = "National Bureau of Economic Research",
 type = "Working Paper",
 series = "Working Paper Series",
 number = "33592",
 year = "2025",
 month = "March",
 doi = {10.3386/w33592},
 URL = "http://www.nber.org/papers/w33592",
 abstract = {Healthcare fraud imposes a sizable cost on U.S. public healthcare budgets and distorts health care provision. We examine the economics of health care fraud and enforcement using theory and data and connect to a growing literature on the topic. We first offer a new economic definition of health care fraud that captures and connects the wide range of activities prosecuted as fraud. We define fraud as any divergence between the care an insurer says a patient qualifies for, the care a provider provides, and the care a provider bills for. Our definition clarifies the economic consequences of different categories of fraud and provides a framework for understanding the slate of existing studies. Next, we examine the incentives for committing and for prosecuting fraud. We show how fraud is driven by a combination of inadequate (expected) penalties for fraud and imperfect reimbursement rates. Public anti-fraud litigation is driven by the relative monetary, political or career returns to prosecuting fraud and by prosecutorial budgets. Finally, we examine the prevalence of health care fraud prosecutions across types of fraud and types of care, and across the US, by machine learning on text data from Department of Justice press releases.},
}

%%%%%%%%%%
%% 2024 %%
%%%%%%%%%%

@article{PughMalani2024,
   author = {Pugh, Sierra and Levin, Andrew T. and Meyerowitz-Katz, Gideon and Soman, Satej and Owusu-Boaitey, Nana and Zwi, Anthony B. and Malani, Anup and Wilson, Ander and Fosdick, Bailey K.},
   title = {A Hierarchical Bayesian Model for Estimating Age-Specific COVID-19 Infection Fatality Rates in Developing Countries},
   journal = {Statistics in Medicine},
   volume = {43},
   number = {30},
   pages = {5667-5680},
   abstract = {ABSTRACT The COVID-19 infection fatality rate (IFR) is the proportion of individuals infected with SARS-CoV-2 who subsequently die. As COVID-19 disproportionately affects older individuals, age-specific IFR estimates are imperative to facilitate comparisons of the impact of COVID-19 between locations and prioritize distribution of scarce resources. However, there lacks a coherent method to synthesize available data to create estimates of IFR and seroprevalence that vary continuously with age and adequately reflect uncertainties inherent in the underlying data. In this article, we introduce a novel Bayesian hierarchical model to estimate IFR as a continuous function of age that acknowledges heterogeneity in population age structure across locations and accounts for uncertainty in the estimates due to seroprevalence sampling variability and the imperfect serology test assays. Our approach simultaneously models test assay characteristics, serology, and death data, where the serology and death data are often available only for binned age groups. Information is shared across locations through hierarchical modeling to improve estimation of the parameters with limited data. Modeling data from 26 developing country locations during the first year of the COVID-19 pandemic, we found seroprevalence did not change dramatically with age, and the IFR at age 60 was above the high-income country estimate for most locations.},
   keywords = {Bayesian inference
hierarchical modeling
prevalence
serology testing
uncertainty quantification},
   ISSN = {0277-6715},
   DOI = {https://doi.org/10.1002/sim.10259},
   url = {https://doi.org/10.1002/sim.10259},
   year = {2024},
   type = {Journal Article}
}

@article{Malani2024covidCellularImmunity,
   author = {Malani, Anup and Aiyar, Jayashree and Sant, Andrea and Kamran, Neha and Mohanan, Manoj and Taneja, Saloni and Woda, Bartek and Zhao, Wanran and Acharya, Anu},
   title = {Comparing population-level humoral and cellular immunity to SARS-Cov-2 in Bangalore, India},
   journal = {Scientific Reports},
   volume = {14},
   number = {1},
   pages = {5758},
   abstract = {Two types of immunity, humoral and cellular, offer protection against COVID. Humoral protection, contributed by circulating neutralizing antibodies, can provide immediate protection but decays more quickly than cellular immunity and can lose effectiveness in the face of mutation and drift in the SARS-CoV-2 spike protein. Therefore, population-level seroprevalence surveys used to estimate population-level immunity may underestimate the degree to which a population is protected against COVID. In early 2021, before India began its vaccination campaign, we tested for humoral and cellular immunity to SARS-Cov-2 in representative samples of slum and non-slum populations in Bangalore, India. We found that 29.7% of samples (unweighted) had IgG antibodies to the spike protein and 15.5% had neutralizing antibodies, but at up to 46% showed evidence of cellular immunity. We also find that prevalence of cellular immunity is significantly higher in slums than in non-slums. These findings suggest (1) that a significantly larger proportion of the population in Bangalore, India, had cellular immunity to SARS-CoV-2 than had humoral immunity, as measured by serological surveys, and (2) that low socio-economic status communities display higher frequency of cellular immunity, likely because of greater exposure to infection due to population density.},
   ISSN = {2045-2322},
   DOI = {10.1038/s41598-024-54922-z},
   url = {https://doi.org/10.1038/s41598-024-54922-z},
   year = {2024},
   type = {Journal Article}
}

@article{SelvavinayagamMalani2024covidSeroTamilNadu,
   author = {Selvavinayagam, T. S. and Somasundaram, Anavarathan and Selvam, Jerard Maria and Sampath, P. and Vijayalakshmi, V. and Kumar, C. Ajith Brabhu and Subramaniam, Sudharshini and Kumarasamy, Parthipan and Raju, S. and Avudaiselvi, R. and Prakash, V. and Yogananth, N. and Subramanian, Gurunathan and Roshini, A. and Dhiliban, D. N. and Imad, Sofia and Tandel, Vaidehi and Parasa, Rajeswari and Sachdeva, Stuti and Ramachandran, Sabareesh and Malani, Anup},
   title = {Contribution of infection and vaccination to population-level seroprevalence through two COVID waves in Tamil Nadu, India},
   journal = {Scientific Reports},
   volume = {14},
   number = {1},
   pages = {2091},
   abstract = {This study employs repeated, large panels of serological surveys to document rapid and substantial waning of SARS-CoV-2 antibodies at the population level and to calculate the extent to which infection and vaccination separately contribute to seroprevalence estimates. Four rounds of serological surveys were conducted, spanning two COVID waves (October 2020 and April–May 2021), in Tamil Nadu (population 72 million) state in India. Each round included representative populations in each district of the state, totaling$\ge$20,000 persons per round. State-level seroprevalence was 31.5\% in round 1 (October–November 2020), after India’s first COVID wave. Seroprevalence fell to 22.9\% in round 2 (April 2021), a roughly one-third decline in 6$>$months, consistent with dramatic waning of SARS-Cov-2 antibodies from natural infection. Seroprevalence rose to 67.1\% by round 3 (June–July 2021), with infections from the Delta-variant induced second COVID wave accounting for 74\% of the increase. Seroprevalence rose to 93.1\% by round 4 (December 2021–January 2022), with vaccinations accounting for 63\% of the increase. Antibodies also appear to wane after vaccination. Seroprevalence in urban areas was higher than in rural areas, but the gap shrunk over time (35.7 v. 25.7\% in round 1, 89.8\% v. 91.4\% in round 4) as the epidemic spread even in low-density rural areas.},
   ISSN = {2045-2322},
   DOI = {10.1038/s41598-023-50338-3},
   url = {https://doi.org/10.1038/s41598-023-50338-3},
   year = {2024},
   type = {Journal Article}
}

% BOOK

@book{Malani2024dekReferenceLegalWeb3,
   author = {Malani, Anup and Henderson, M. Todd},
   title = {Legal Matters in Web 3: A Desk Reference},
   url = {https://www.law.uchicago.edu/sites/default/files/2024-09/LM%20Web3%20Desk%20Reference%20-%20Formatted%20-%20240907.pdf},
   year = {2024},
   type = {Book}
}

% WORKING-PAPER

@article{Malani2024pricingHealthInsuranceIndia,
   author = {Malani, Anup and Kinnan, Cynthia and Conti, Gabriella and Imai, Kosuke and Miller, Morgen and Swaminathan, Shailender and Voena, Alessandra and Woda, Bartosz},
   title = {Evaluating and Pricing Health Insurance in Lower-income Countries: A Field Experiment in India},
   journal = {National Bureau of Economic Research Working Paper Series},
   volume = {No. 32239},
   note = {Author contact info: Anup Malani University of Chicago Law School 1111 E. 60th Street Chicago, IL 60637 Tel: 773/702-9602 E-Mail: amalani@uchicago.edu Cynthia Kinnan Department of Economics Tufts University 177 College Ave #620 Medford, MA 02155 E-Mail: cynthia.kinnan@tufts.edu Gabriella Conti Department of Economics University College London Drayton House 30 Gordon Street London WC1H 0AX United Kingdom Tel: Tel. +44(0)20 76794 1731 E-Mail: gabriella.conti@ucl.ac.uk Kosuke Imai Harvard University 1737 Cambridge Street Cambridge, MA 02138 E-Mail: imai@harvard.edu Morgen Miller University of Chicago Law School 1111 E. 60th St. Chicago, IL 60637 E-Mail: mmmiller@uchicago.edu Shailender Swaminathan Sai University SSPDL Building, Beta Block 4th Floor Navalur Old Mahabalipuram Rd Chennai India E-Mail: shailender.s@saiuniversity.edu.in Alessandra Voena Stanford University Landau Economics Building 579 Jane Stanford Way Stanford, CA 94305 Tel: 650/498-1096 E-Mail: avoena@stanford.edu Bartosz Woda University of Chicago Law School 1111 E. 60th St. Chicago, IL 60637 USA E-Mail: bartekwoda@gmail.com},
   abstract = {Universal health coverage is a widely shared goal across lower-income countries. We conducted a large-scale, 4-year trial that randomized premiums and subsidies for India’s first national, public hospital insurance program, called RSBY. We find substantial demand (∼ 60\% uptake) even when consumers were charged a price equal to the premium the government paid for insurance. We also find substantial adverse selection into insurance at positive prices. Insurance enrollment increases insurance utilization, partly due to spillovers from use of insurance by neighbors. However, healthcare utilization does not rise substantially, suggesting the primary benefit of insurance is financial. Many enrollees attempted to use insurance but failed, suggesting that learning is critical to the success of public insurance. We find very few statistically significant impacts of insurance access or enrollment on health. Because there is substantial willingness-to- pay for insurance, and given how distortionary it is to raise revenue in the Indian context, we calculate that our sample population should be charged a premium for RSBY between 67-95\% of average costs (INR 528-1052, $30-60) rather than a zero premium to maximize the marginal value of public funds.},
   DOI = {10.3386/w32239},
   url = {http://www.nber.org/papers/w32239},
   year = {2024},
   type = {Journal Article}
}

@article{RosenbaumMalani2024earlyLifeEffects,
   author = {Rosenbaum, Stacy and Malani, Anup and Lea, Amanda J. and Tung, Jenny and Alberts, Susan C. and Archie, Elizabeth A.},
   title = {Testing frameworks for early life effects: the developmental constraints and adaptive response hypotheses do not explain key fertility outcomes in wild female baboons},
   journal = {bioRxiv},
   pages = {2024.04.23.590627},
   abstract = {In evolutionary ecology, two classes of explanations are frequently invoked to explain “early life effects” on adult outcomes. Developmental constraints (DC) explanations contend that costs of early adversity arise from limitations adversity places on optimal development. Adaptive response (AR) hypotheses propose that later life outcomes will be worse when early and adult environments are poorly “matched.” Here, we use recently proposed mathematical definitions for these hypotheses and a quadratic-regression based approach to test the long-term consequences of variation in developmental environments on fertility in wild baboons. We evaluate whether low rainfall and/or dominance rank during development predict three female fertility measures in adulthood, and whether any observed relationships are consistent with DC and/or AR. Neither rainfall during development nor the difference between rainfall in development and adulthood predicted any fertility measures. Females who were low-ranking during development had an elevated risk of losing infants later in life, and greater change in rank between development and adulthood predicted greater risk of infant loss. However, both effects were statistically marginal and consistent with alternative explanations, including adult environmental quality effects. Consequently, our data do not provide compelling support for either of these common explanations for the evolution of early life effects.Competing Interest Statement: The authors have declared no competing interest.},
   DOI = {10.1101/2024.04.23.590627},
   url = {http://biorxiv.org/content/early/2024/04/28/2024.04.23.590627.abstract},
   year = {2024},
   type = {Journal Article}
}

@article{JaffeMalani2024creditReducesValueInsurance,
   author = {Jaffe, Sonia and Malani, Anup and Reif, Julian},
   title = {Access to Credit Reduces the Value of Insurance},
   journal = {National Bureau of Economic Research Working Paper Series},
   volume = {No. 32395},
   note = {Author contact info: Sonia Jaffe Microsoft Principal Researcher E-Mail: spj@uchicago.edu Anup Malani University of Chicago Law School 1111 E. 60th Street Chicago, IL 60637 Tel: 773/702-9602 E-Mail: amalani@uchicago.edu Julian Reif Department of Finance University of Illinois Urbana-Champaign 515 E. Gregory Street Urbana, IL 61820 Tel: 217/300-0169 E-Mail: jreif@illinois.edu},
   abstract = {We analyze the value of insurance when individuals have access to credit markets. Loans allow consumers to smooth financial shocks over time, decreasing the value of consumption smoothing from insurance. We derive formulas for the value of insurance that can be taken to data, and show how that value depends on individual characteristics and features of loans. We estimate that access to a five-year loan decreases the values of community- and experience-rated insurance for the average beneficiary by \$232-\$366 (58--61\%). Even for the sickest decile, this loan access reduces the value of community-rated insurance by \$1,099 (17\%).},
   DOI = {10.3386/w32395},
   url = {http://www.nber.org/papers/w32395},
   year = {2024},
   type = {Journal Article}
}

@article{AbbotMalani2024limitsGrowth,
   author = {Abbot, Dorian S. and Malani, Anup},
   title = {Revisiting the Physical Limits to Economic Growth },
   journal = {Working Paper},
   year = {2024},
   type = {Journal Article}
}

@article{RN6625,
   author = {Holden, Richard T. and Malani, Anup and Teh, Chris},
   title = {Allocating Scarce Information},
   journal = {National Bureau of Economic Research Working Paper Series},
   volume = {No. 29846},
   note = {Author contact info: Richard T. Holden University of New South Wales Room 470B Sydney, NSW, 2052, AUSTRALIA Australia Tel: +61 2 9385 4700 Fax: +1 773 409 5383 E-Mail: richard.holden@unsw.edu.au Anup Malani University of Chicago Law School 1111 E. 60th Street Chicago, IL 60637 Tel: 773/702-9602 E-Mail: amalani@uchicago.edu Chris Teh University of New South Wales Sydney, NSW 2052 Australia E-Mail: chris.teh@unsw.edu.au},
   abstract = {Sender conveys scarce information to a number of receivers to maximize the sum of receiver payoffs. Each receiver’s payoff depends on the state of the world and an action she takes. The optimal action is state contingent. Under mild regularity conditions, we show that the payoff of each receiver is convex in the amount of information she receives. Thus, it is optimal for Sender to target information to a single receiver. We then study four extensions in which interior information allocations are optimal.},
   DOI = {10.3386/w29846},
   url = {http://www.nber.org/papers/w29846},
   year = {2022},
   type = {Journal Article}
}

@article{RN6326,
   author = {Malani, Anup and Zhao, Wanran},
   title = {The Mortality Burden from COVID in Low-Income Settings: Evidence from Verbal Autopsies in India},
   journal = {medRxiv},
   pages = {2024.01.02.24300728},
   abstract = {Measuring the mortality burden of SARS-CoV-2 in lower-income countries is difficult because death registries are incomplete and lack cause of death. We address this problem in India, which had the second-highest number of officially reported cases. We completed WHO-compliant verbal autopsy (VA) surveys on roughly 20,000 deaths drawn from a population-representative sample. SARS-CoV-2 deaths spike in June 2020, just after India’s lockdown, and in May 2021, after its second wave. During those spikes the virus is responsible for 23.3\% and 35.8\% of all deaths, respectively. Cardiovascular deaths also spike during the start of the pandemic. We find that the death rate rises by 81\% during the pandemic, SARS-CoV-2 is responsible for 33\% of these excess deaths, and cardiovascular disease for 23\% of these deaths.Competing Interest StatementThe authors have declared no competing interest.Funding StatementThe Becker-Friedman Institute at the University of Chicago provided a site license to the Consumer Pyramids Household Survey data (from the Centre for Monitoring Indian Economy) to which we linked the verbal autopsy data. The verbal autopsy data were made available without charge by the Centre for Monitoring Indian Economy.Author DeclarationsI confirm all relevant ethical guidelines have been followed, and any necessary IRB and/or ethics committee approvals have been obtained.YesThe details of the IRB/oversight body that provided approval or exemption for the research described are given below:The University of Chicago IRB determined this study (Case No. IRB22-0119) was not human subjects research as no identifiable information is shared and the data are or will be publicly accessible.I confirm that all necessary patient/participant consent has been obtained and the appropriate institutional forms have been archived, and that any patient/participant/sample identifiers included were not known to anyone (e.g., hospital staff, patients or participants themselves) outside the research group so cannot be used to identify individuals.YesI understand that all clinical trials and any other prospective interventional studies must be registered with an ICMJE-approved registry, such as ClinicalTrials.gov. I confirm that any such study reported in the manuscript has been registered and the trial registration ID is provided (note: if posting a prospective study registered retrospectively, please provide a statement in the trial ID field explaining why the study was not registered in advance).YesI have followed all appropriate research reporting guidelines, such as any relevant EQUATOR Network research reporting checklist(s) and other pertinent material, if applicable.YesAll data in the present study will be made available to the public via the open data website openmortality.org. https://consumerpyramidsdx.cmie.com/},
   DOI = {10.1101/2024.01.02.24300728},
   url = {https://www.medrxiv.org/content/medrxiv/early/2024/01/03/2024.01.02.24300728.full.pdf},
   year = {2024},
   type = {Journal Article}
}



%%%%%%%%%%
%% 2023 %%
%%%%%%%%%%

@article{Malani2023devConstraintsPARModel,
   author = {Malani, Anup and Archie, Elizabeth A. and Rosenbaum, Stacy},
   title = {Conceptual and analytical approaches for modelling the developmental origins of inequality},
   journal = {Philosophical Transactions of the Royal Society B: Biological Sciences},
   volume = {378},
   number = {1883},
   pages = {20220306},
   note = {doi: 10.1098/rstb.2022.0306},
   DOI = {10.1098/rstb.2022.0306},
   url = {https://doi.org/10.1098/rstb.2022.0306},
   year = {2023},
   type = {Journal Article}
}

@article{Malani2023covidLessons,
   author = {Malani, Anup},
   title = {Lessons from Disease and Economic Surveillance during COVID in India},
   journal = {India Policy Forum},
   volume = {19},
   number = {1},
   pages = {121-180},
   year = {2023},
   type = {Journal Article}
}


@article{DuMalani2023fractionalDosing,
   author = {Du, Zhanwei and Wang, Lin and Bai, Yuan and Feng, Shuo and Ramachandran, Sabareesh and Lim, Wey Wen and Lau, Eric H. Y. and Malani, Anup and Cowling, Benjamin J.},
   title = {Cost effectiveness of fractional doses of COVID-19 vaccine boosters in India},
   journal = {Med},
   volume = {4},
   number = {3},
   pages = {182-190.e3},
   abstract = {Summary Background Coronavirus disease 2019 (COVID-19) continues to be a major global public health crisis that exacts significant human and economic costs. Booster vaccination of individuals can improve waning immunity and reduce the impact of community epidemics. Methods Using an epidemiological model that incorporates population-level severe acute respiratory syndrome coronavirus 2 (SARS-CoV-2) transmission and waning of vaccine-derived immunity, we identify the hypothetical potential of mass vaccination with fractionated vaccine doses specific to ChAdOx1 nCoV-19 (AZD1222 [Covishield]; AstraZeneca) as an optimal and cost-effective strategy in India’s Omicron outbreak. Findings We find that the optimal strategy is 1/8 fractional dosing under mild (Re ∼ 1.2) and rapid (Re ∼ 5) transmission scenarios, leading to an estimated $6 (95\% confidence interval [CI]: −13, 26) billion and $2 (95\% CI: −26, 30) billion in health-related net monetary benefit over 200 days, respectively. Rapid and broad use of fractional dosing for boosters, together with delivery costs divided by fractionation, could substantially gain more net monetary benefit by $11 (95\% CI: −10, 33) and $2 (95\% CI: −23, 28) billion, respectively, under the mild and rapid transmission scenarios. Conclusions Mass vaccination with fractional doses of COVID-19 vaccines to boost immunity in a vaccinated population could be a cost-effective strategy for mitigating the public health costs of resurgences caused by vaccine-evasive variants, and fractional dosing deserves further clinical and regulatory evaluation. Funding Financial support was provided by the AIR@InnoHK Program from Innovation and Technology Commission of the Government of the Hong Kong Special Administrative Region.},
   keywords = {COVID-19
SARS-CoV-2
dose fractionation
vaccine booster
epidemiological model
cost effectiveness},
   ISSN = {2666-6340},
   DOI = {https://doi.org/10.1016/j.medj.2023.02.001},
   url = {https://www.sciencedirect.com/science/article/pii/S2666634023000612},
   year = {2023},
   type = {Journal Article}
}

@article{ShafrinMalani2023drugPriceNegotsIRA,
   author = {Shafrin, Jason and Lakdawalla, Darius N and Doshi, Jalpa A and Garrison Jr, Louis P and Malani, Anup and Neumann, Peter J and Phelps, Charles E and Towse, Adrian and Willke, Richard J},
   title = {A strategy for value-based drug pricing under the inflation reduction act},
   journal = {Health Affairs Forefront},
   DOI = {10.1377/forefront.20230503.153705},
   year = {2023},
   type = {Journal Article}
}

@article{JiangMalani2023powerSpilloverTwoStageExperiment,
   author = {Jiang, Zhichao and Imai, Kosuke and Malani, Anup},
   title = {Statistical inference and power analysis for direct and spillover effects in two-stage randomized experiments},
   journal = {Biometrics},
   volume = {79},
   number = {3},
   pages = {2370-2381},
   abstract = {Abstract Two-stage randomized experiments become an increasingly popular experimental design for causal inference when the outcome of one unit may be affected by the treatment assignments of other units in the same cluster. In this paper, we provide a methodological framework for general tools of statistical inference and power analysis for two-stage randomized experiments. Under the randomization-based framework, we consider the estimation of a new direct effect of interest as well as the average direct and spillover effects studied in the literature. We provide unbiased estimators of these causal quantities and their conservative variance estimators in a general setting. Using these results, we then develop hypothesis testing procedures and derive sample size formulas. We theoretically compare the two-stage randomized design with the completely randomized and cluster randomized designs, which represent two limiting designs. Finally, we conduct simulation studies to evaluate the empirical performance of our sample size formulas. For empirical illustration, the proposed methodology is applied to the randomized evaluation of the Indian National Health Insurance Program. An open-source software package is available for implementing the proposed methodology.},
   keywords = {experimental design
interference between units
partial interference
spillover effects
statistical power},
   ISSN = {0006-341X},
   DOI = {https://doi.org/10.1111/biom.13782},
   url = {https://doi.org/10.1111/biom.13782},
   year = {2023},
   type = {Journal Article}
}

% NON-JOURNAL

@book{CovidCrisisGroup2023lessonsCovidWar,
   author = {Covid Crisis Group,},
   title = {Lessons from the Covid War: An Investigative Report},
   publisher = {PublicAffairs},
   ISBN = {9781541703810},
   url = {https://books.google.com/books?id=gEedEAAAQBAJ},
   year = {2023},
   type = {Book}
}


%%%%%%%%%%
%% 2022 %%
%%%%%%%%%%

@article{LevinMalani2023covidIFRdevelopingCountries,
   author = {Levin, Andrew T. and Owusu-Boaitey, Nana and Pugh, Sierra and Fosdick, Bailey K. and Zwi, Anthony B. and Malani, Anup and Soman, Satej and Besançon, Lonni and Kashnitsky, Ilya and Ganesh, Sachin and McLaughlin, Aloysius and Song, Gayeong and Uhm, Rine and Herrera-Esposito, Daniel and de los Campos, Gustavo and Peçanha Antonio, Ana Carolina and Tadese, Enyew Birru and Meyerowitz-Katz, Gideon},
   title = {Assessing the burden of COVID-19 in developing countries: systematic review, meta-analysis and public policy implications},
   journal = {BMJ Global Health},
   volume = {7},
   number = {5},
   pages = {e008477},
   abstract = {Introduction The infection fatality rate (IFR) of COVID-19 has been carefully measured and analysed in high-income countries, whereas there has been no systematic analysis of age-specific seroprevalence or IFR for developing countries.Methods We systematically reviewed the literature to identify all COVID-19 serology studies in developing countries that were conducted using representative samples collected by February 2021. For each of the antibody assays used in these serology studies, we identified data on assay characteristics, including the extent of seroreversion over time. We analysed the serology data using a Bayesian model that incorporates conventional sampling uncertainty as well as uncertainties about assay sensitivity and specificity. We then calculated IFRs using individual case reports or aggregated public health updates, including age-specific estimates whenever feasible.Results In most locations in developing countries, seroprevalence among older adults was similar to that of younger age cohorts, underscoring the limited capacity that these nations have to protect older age groups.Age-specific IFRs were roughly 2 times higher than in high-income countries. The median value of the population IFR was about 0.5%, similar to that of high-income countries, because disparities in healthcare access were roughly offset by differences in population age structure.Conclusion The burden of COVID-19 is far higher in developing countries than in high-income countries, reflecting a combination of elevated transmission to middle-aged and older adults as well as limited access to adequate healthcare. These results underscore the critical need to ensure medical equity to populations in developing countries through provision of vaccine doses and effective medications.Data are available in a public, open access repository.},
   DOI = {10.1136/bmjgh-2022-008477},
   url = {http://gh.bmj.com/content/7/5/e008477.abstract},
   year = {2022},
   type = {Journal Article}
}

@article{BakerMalani2022accuracy,
   author = {Baker, Scott and Malani, Anup},
   title = {Compromising Accuracy to Encourage Regulatory Participation},
   journal = {The Journal of Legal Studies},
   volume = {51},
   number = {1},
   pages = {1-38},
   note = {doi: 10.1086/719238},
   abstract = {AbstractThis paper examines the value of accuracy in voluntary or opt-in regulatory regimes. We show that if welfare depends primarily on the ability to identify noncompliant firms, tolerating mistakes in concluding that firms meet regulatory standards (false positives) can improve welfare. Consumers or investors anticipate the regulatory error rate and discount the positive message of a compliance finding. Welfare is nonetheless improved because the mistaken exoneration acts as an incentive for noncompliant firms to submit to regulatory scrutiny, and as a result some noncompliant firms are unmasked.},
   ISSN = {0047-2530},
   DOI = {10.1086/719238},
   url = {https://doi.org/10.1086/719238},
   year = {2022},
   type = {Journal Article}
}

@article{Malani2022covidMortalityIndia,
   author = {Malani, Anup and Ramachandran, Sabareesh},
   title = {Using household rosters from survey data to estimate all-cause excess death rates during the COVID pandemic in India},
   journal = {Journal of Development Economics},
   volume = {159},
   pages = {102988},
   abstract = {Official statistics on deaths from COVID undercount deaths due to lack of testing. In developed countries, death registries are used to estimate excess deaths due to COVID during the pandemic. However, few developing countries had complete death registries even before the pandemic and the pandemic further stressed administrative capacities. As a substitute, we estimate all-cause excess deaths in India using the member rosters of a large, representative household panel survey. We estimate roughly 4.2 million excess deaths during the pandemic through February 2022. We cannot demonstrate causality between COVID and deaths, but the timing and age structure of deaths is consistent with the COVID pandemic and excess deaths are positively correlated with reported infections. Finally, we find that excess deaths were higher among higher-income persons and were negatively associated with mobility. The methods in this paper can be used in countries with a household panel to measure health-related demographic indicators.},
   keywords = {SARS-Cov-2
COVID
Pandemic
India
Mortality rate
Excess deaths
Household survey},
   ISSN = {0304-3878},
   DOI = {https://doi.org/10.1016/j.jdeveco.2022.102988},
   url = {https://www.sciencedirect.com/science/article/pii/S0304387822001304},
   year = {2022},
   type = {Journal Article}
}

@article{DuMalani2022modelingFractionation,
   author = {Du, Zhanwei and Wang, Lin and Pandey, Abhishek and Lim, Wey Wen and Chinazzi, Matteo and Piontti, Ana Pastore y and Lau, Eric H. Y. and Wu, Peng and Malani, Anup and Cobey, Sarah and Cowling, Benjamin J.},
   title = {Modeling comparative cost-effectiveness of SARS-CoV-2 vaccine dose fractionation in India},
   journal = {Nature Medicine},
   volume = {28},
   number = {5},
   pages = {934-938},
   abstract = {Given global Coronavirus Disease 2019 (COVID-19) vaccine shortages and inequity of vaccine distributions, fractionation of vaccine doses might be an effective strategy for reducing public health and economic burden, notwithstanding the emergence of new variants of concern. In this study, we developed a multi-scale model incorporating population-level transmission and individual-level vaccination to estimate the costs of hospitalization and vaccination and the economic benefits of reducing COVID-19 deaths due to dose-fractionation strategies in India. We used large-scale survey data of the willingness to pay together with data of vaccine and hospital admission costs to build the model. We found that fractional doses of vaccines could be an economically viable vaccination strategy compared to alternatives of either full-dose vaccination or no vaccination. Dose-sparing strategies could save a large number of lives, even with the emergence of new variants with higher transmissibility.},
   ISSN = {1546-170X},
   DOI = {10.1038/s41591-022-01736-z},
   url = {https://doi.org/10.1038/s41591-022-01736-z},
   year = {2022},
   type = {Journal Article}
}

@article{HoldenMalani2022blockchainVelocityICO,
   author = {Holden, Richard and Malani, Anup},
   title = {An Examination of Velocity and Initial Coin Offerings},
   journal = {Management Science},
   volume = {68},
   number = {12},
   pages = {9026-9041},
   note = {doi: 10.1287/mnsc.2022.4314},
   abstract = {Blockchain technology offers firms a novel method of raising capital via so-called initial coin offerings (ICOs). In the most common form of an ICO, a firm creates digital assets called ?utility tokens? that are tracked on a blockchain-based ledger, requires that its product be purchased only with those tokens, and then, raises capital by selling these tokens to investors prior to creating any saleable product. (Some nonfungible tokens (NFTs) may function in a similar fashion.) We model a fundamental paradox with the use of ICOs involving utility tokens and similar structures. To increase capital raised by an ICO, the firm may attempt to reduce blockchain operating costs, thus expanding the quantity of goods sold. However, because of the mechanics of miner compensation, doing so increases the number of utility token transactions that take place in any time interval (i.e., increases token velocity and thus, the effective supply of tokens). By Fisher?s equation, this lowers the dollar value of tokens and the amount investors are willing to pay for them. We show that this paradox limits the value of utility token ICOs as an alternative to traditional financing options. We discuss alternatives to and variations of utility tokens that can mitigate the conundrum and promote ICOs as a more viable form of financing.This paper was accepted by Joshua Gans, business strategy.},
   ISSN = {0025-1909},
   DOI = {10.1287/mnsc.2022.4314},
   url = {https://doi.org/10.1287/mnsc.2022.4314},
   year = {2022},
   type = {Journal Article}
}

@article{HoldenMalani2022lawEconBlockchain,
   author = {Holden, Richard and Malani, Anup},
   title = {The Law and Economics of Blockchain},
   journal = {Annual Review of Law and Social Science},
   volume = {18},
   pages = {297-313},
   abstract = {This article examines the implications of Distributed Ledger Technology (a.k.a. blockchain) for several areas of law. While cryptocurrencies have received much attention, the implications of DLT are potentially far reaching. DLT raises interesting and important questions relating to rules of evidence, surrounding issues like hearsay and authentication. The advent of initial coin offerings has implications not only for how firms are financed but also for securities law in regulating such offerings. Cryptocurrencies themselves (e.g., Bitcoin) have raised serious issues for tax avoidance and taxation law. Relatedly, the rise of cryptocurrencies raises issues regarding the relationship between private and nationally issued currencies, and even the role and efficacy of monetary policy. Finally, DLT has practical implications for election law and voter turnout, with such technology already begun to be used for security purposes in online voting and permitted in 32 US states.},
   keywords = {regulation
contracts
blockchain
law
distributed ledger},
   ISSN = {1550-3631},
   DOI = {https://doi.org/10.1146/annurev-lawsocsci-011921-060322},
   url = {https://www.annualreviews.org/content/journals/10.1146/annurev-lawsocsci-011921-060322},
   year = {2022},
   type = {Journal Article}
}

@article{WenMalani2022flueEvolutionVaccine,
   author = {Wen, Frank T. and Malani, Anup and Cobey, Sarah},
   title = {The Potential Beneficial Effects of Vaccination on Antigenically Evolving Pathogens},
   journal = {The American Naturalist},
   volume = {199},
   number = {2},
   pages = {223-237},
   note = {doi: 10.1086/717410},
   abstract = {AbstractAlthough vaccines against antigenically evolving pathogens such as seasonal influenza ; and are designed to protect against circulating strains by affecting the emergence and transmission of antigenically divergent strains, they might in theory also be able to change the rate of antigenic evolution. Vaccination might slow antigenic evolution by increasing immunity, reducing the overall prevalence or population size of the pathogen. This reduction could decrease the supply and growth rates of mutants and might thereby slow adaptation. But vaccination might accelerate antigenic evolution by increasing the transmission advantage of more antigenically diverged strains relative to less diverged strains (i.e., by positive selection). Such evolutionary effects could affect vaccination?s direct benefits to individuals and indirect benefits to the host population (i.e., the private and social benefits). To investigate these potential impacts, we simulated vaccination against a continuously circulating influenza-like pathogen in a simple population. On average, more vaccination decreased the incidence of infection. Notably, this decrease was driven partly by a vaccine-induced decline in the rate of antigenic evolution. To understand how the evolutionary effects of vaccines might affect their social and private benefits, we fitted linear panel models to simulated data. By slowing evolution, vaccination increased the social benefit and decreased the private benefit. Thus, vaccination?s potential social and private benefits may differ from current theory, which omits evolutionary effects. These results suggest that conventional vaccines against influenza and other antigenically evolving pathogens, if protective against transmission and given to the appropriate populations, could further reduce disease burden by slowing antigenic evolution.},
   ISSN = {0003-0147},
   DOI = {10.1086/717410},
   url = {https://doi.org/10.1086/717410},
   year = {2021},
   type = {Journal Article}
}

@article{ShengMalani2022slumMobilityCOVID,
   author = {Sheng, Jaymee and Malani, Anup and Goel, Ashish and Botla, Purushotham},
   title = {JUE insights: Does mobility explain why slums were hit harder by COVID-19 in Mumbai, India?},
   journal = {Journal of Urban Economics},
   volume = {127},
   pages = {103357},
   abstract = {SARS-CoV-2 has had a greater burden, as measured by rate of infection, in poorer communities within cities. For example, 55% of Mumbai slums residents had antibodies to COVID-19, 3.2 times the seroprevalence in non-slum areas of the city according to a sero-survey done in July 2020. One explanation is that government suppression was less severe in poorer communities, either because the poor were more likely to be exempt or unable to comply. Another explanation is that effective suppression itself accelerated the epidemic in poor neighborhoods because households are more crowded and residents share toilet and water facilities. We show there is little evidence for the first hypothesis in the context of Mumbai. Using location data from smart phones, we find that slum residents had nominally but not significantly (economically or statistically) higher mobility than non-slums prior to the sero-survey. We also find little evidence that mobility in non-slums was lower than in slums during lockdown, a subset of the period before the survey.},
   keywords = {COVID-19
Infection
Epidemic
Slums
Mobility
Inequality},
   ISSN = {0094-1190},
   DOI = {https://doi.org/10.1016/j.jue.2021.103357},
   url = {https://www.sciencedirect.com/science/article/pii/S0094119021000395},
   year = {2022},
   type = {Journal Article}
}

@article{TandelMalani2022slumCOVIDInfra,
   author = {Tandel, Vaidehi and Gandhi, Sahil and Patranabis, Shaonlee and Bettencourt, Luís M. A. and Malani, Anup},
   title = {Infrastructure, enforcement, and COVID-19 in Mumbai slums: A first look},
   journal = {Journal of Regional Science},
   volume = {62},
   number = {3},
   pages = {645-669},
   abstract = {Abstract This study is among the first to investigate whether patterns of access to basic services could explain the disproportionately severe impact of COVID-19 in slums. Using geolocated containment zones and COVID-19 case data for Mumbai, India's most populous city, we find that cases and case fatality rates are higher in slums compared with formal residential buildings. Our results show that access to toilets for men is associated with lower COVID-19 prevalence. However, the effect is opposite in the case of toilets for women. This could be because limited hours for safely using toilets and higher waiting times increase the risk of exposure, and women and children sharing toilet facilities results in crowding. Proximity to water pipelines has no effect on prevalence, likely because slumdwellers are disconnected from formal water supply networks. Indoor crowding does not seem to have an effect on case prevalence. Finally, while police capacity?measured by number of police station outposts?is associated with lower prevalence in nonslum areas, indicating effective enforcement of containment, this relationship does not hold in slums. The study highlights the urgency of finding viable solutions for slum improvement and upgrading to mitigate the effects of contagion for some of the most vulnerable populations.},
   keywords = {basic services
COVID-19
India
Mumbai
slums
spatial inequality},
   ISSN = {0022-4146},
   DOI = {https://doi.org/10.1111/jors.12552},
   url = {https://doi.org/10.1111/jors.12552},
   year = {2022},
   type = {Journal Article}
}

% NON-JOURNAL


@inbook{MohananMalani2022seroPolicymakingCovid,
   author = {Mohanan, Manoj and Malani, Anup and Acharya, Anu},
   title = {The Role of Seroprevalence in Evidence for Policy-Making During COVID-19},
   booktitle = {Flattening the Curve},
   series = {World Scientific Series in Grand Public Policy Challenges of the 21st Century},
   publisher = {WORLD SCIENTIFIC},
   volume = {Volume 4},
   pages = {187-205},
   note = {doi:10.1142/9789811262739_0007},
   abstract = {Seroprevalence studies estimate the share of population that has been infected by a virus and can provide valuable information for policymakers as they respond to a pandemic, especially in its early stages. The spread of COVID-19 globally, as measured by seroprevalence studies, has been considerably more extensive than what is perceived from the case numbers reported by governments worldwide. The inability to incorporate evidence from such studies in a timely manner contributed to missed opportunities in the policy response to COVID-19. In this chapter, we discuss the potential role that seroprevalence studies can play in providing timely evidence for policy-making as well as the factors that lead to impediments in the adoption of such evidence. Building on our own experience with conducting some of the earliest seroprevalence studies in India, we describe the challenges of policy-making in crises and of incorporating uncertainty into this process, and we also examine some of the technical challenges in conducting these studies. In addition, we discuss policy responses that could be informed by evidence from such studies and highlight the lessons learned that might inform efforts in a future crisis, should it ever occur.},
   ISBN = {978-981-12-6272-2},
   DOI = {doi:10.1142/9789811262739_0007
10.1142/9789811262739_0007},
   url = {https://doi.org/10.1142/9789811262739_0007},
   year = {2022},
   type = {Book Section}
}



% WORKING-PAPER

@article{Malani2022vaccineAllocationDemand,
   author = {Malani, Anup and Soman, Satej and Ramachandran, Sabareesh and Chen, Alice and Lakdawalla, Darius N.},
   title = {Vaccine Allocation Priorities Using Disease Surveillance and Economic Data},
   journal = {National Bureau of Economic Research Working Paper Series},
   volume = {No. 29682},
   note = {Author contact info: Anup Malani University of Chicago Law School 1111 E. 60th Street Chicago, IL 60637 Tel: 773/702-9602 E-Mail: amalani@uchicago.edu Satej Soman School of Information},
   abstract = {Vaccination is a critical tool, along with suppression and treatment, for controlling epidemics such as SARS-CoV-2. To maximize the impact of vaccination, doses should be allocated to the highest value targets, accounting for health and potential economic benefits. We examine what allocation strategy is optimal and how to translate that strategy into actionable procurement decisions in the context of India. We compare 3 different allocation strategies (oldest first, highest contact rate first, random order) across 4 outcomes (lives saved, life-years saved, value of statistical lives saved, value of statistical life-years saved). We make 3 methodological contributions. First, we estimate the incremental health benefit of vaccination using novel, local seroprevalence data from India. Second, we estimate the value of statistical life-years using disaggregated, monthly data on consumption during the pandemic. Third, and most importantly, we estimate social demand curves for vaccines that can practically guide government procurement decisions. Our analysis yields 4 novel findings. First, the need to speed-up vaccination does not justify deviation from elderly-first prioritization. Second, much of the value of vaccination comes from improvements in consumption rather than longevity. Moreover, vaccination increases the value of a life year because it increases consumption. Third, social demand for vaccination falls over time as natural immunity from infections increases. Therefore, the slower a country vaccinates its population, the fewer doses it should procure. Fourth, there is enough variation in consumption and infection risk that it makes sense to vaccinate some areas before others. Our approach of connecting epidemiological models and data on health and consumption to economic valuation methods generalizes to other infection control strategies, such as suppression, and public health crises, such as influenza and HIV.},
   DOI = {10.3386/w29682},
   url = {http://www.nber.org/papers/w29682},
   year = {2022},
   type = {Journal Article}
}



%%%%%%%%%%
%% 2021 %%
%%%%%%%%%%

@article{Malani2021mortaPeril,
   author = {Malani, Anup},
   title = {Still in Mortal Peril? Recent Research Suggests a New Agenda for Health Care Reform},
   journal = {The Journal of Legal Studies},
   volume = {50},
   number = {S2},
   pages = {S239-S273},
   note = {doi: 10.1086/708094},
   abstract = {AbstractThe Affordable Care Act expanded funding for health insurance and regulated health insurance markets. These actions followed from two premises. First, nearly 50 million people lacked health insurance in 2010 and were unable to obtain critical medical care. Second, the primary threat to health insurance markets is adverse selection: because healthy individuals are reluctant to join insurance pools, they must be mandated to purchase insurance. Research on those premises yields inconsistent conclusions. I review this research and offer an alternative framework for understanding health-related risk and health care financing in the United States. Three implications follow. First, the best way to tackle health-related risk is to encourage innovation and eliminate inefficiencies that inflate health care prices. Second, insurance markets should be stabilized by allowing experience rating and long-term contracts instead of requiring community rating and imposing insurance mandates. Third, redistribution is better accomplished via targeted premium subsidies than insurance-pricing regulations.},
   ISSN = {0047-2530},
   DOI = {10.1086/708094},
   url = {https://doi.org/10.1086/708094},
   year = {2021},
   type = {Journal Article}
}

@article{CaiMalani2021covidIFRIndia,
   author = {Cai, Rebecca and Novosad, Paul and Tandel, Vaidehi and Asher, Sam and Malani, Anup},
   title = {Representative estimates of COVID-19 infection fatality rates from four locations in India: cross-sectional study},
   journal = {BMJ Open},
   volume = {11},
   number = {10},
   pages = {e050920},
   abstract = {Objectives To estimate age-specific and sex-specific mortality risk among all SARS-CoV-2 infections in four settings in India, a major lower-middle-income country and to compare age trends in mortality with similar estimates in high-income countries.Design Cross-sectional study.Setting India, multiple regions representing combined population &amp;gt;150 million.Participants Aggregate infection counts were drawn from four large population-representative prevalence/seroprevalence surveys. Data on corresponding number of deaths were drawn from official government reports of confirmed SARS-CoV-2 deaths.Primary and secondary outcome measures The primary outcome was age-specific and sex-specific infection fatality rate (IFR), estimated as the number of confirmed deaths per infection. The secondary outcome was the slope of the IFR-by-age function, representing increased risk associated with age.Results Among males aged 50–89, measured IFR was 0.12\% in Karnataka (95\% CI 0.09\% to 0.15\%), 0.42\% in Tamil Nadu (95\% CI 0.39\% to 0.45\%), 0.53\% in Mumbai (95\% CI 0.52\% to 0.54\%) and an imprecise 5.64\% (95\% CI 0\% to 11.16\%) among migrants returning to Bihar. Estimated IFR was approximately twice as high for males as for females, heterogeneous across contexts and rose less dramatically at older ages compared with similar studies in high-income countries.Conclusions Estimated age-specific IFRs during the first wave varied substantially across India. While estimated IFRs in Mumbai, Karnataka and Tamil Nadu were considerably lower than comparable estimates from high-income countries, adjustment for under-reporting based on crude estimates of excess mortality puts them almost exactly equal with higher-income country benchmarks. In a marginalised migrant population, estimated IFRs were much higher than in other contexts around the world. Estimated IFRs suggest that the elderly in India are at an advantage relative to peers in high-income countries. Our findings suggest that the standard estimation approach may substantially underestimate IFR in low-income settings due to under-reporting of COVID-19 deaths, and that COVID-19 IFRs may be similar in low-income and high-income settings.Data are available in a public, open access repository. Replication code, data dictionary, and data will be posted in a public repository on Github. The repository will include all data on demographics and COVID-19 deaths by location, seroprevalence aggregates for Mumbai, Karnataka, and Tamil Nadu, and mortality rates by age and gender for migrants from Bihar. We do not have permission to share seroprevalence microdata. Replication code will be provided to reconstruct all results in the paper from these data.},
   DOI = {10.1136/bmjopen-2021-050920},
   url = {http://bmjopen.bmj.com/content/11/10/e050920.abstract},
   year = {2021},
   type = {Journal Article}
}

@article{MohananMalani2021covidSeroKarnatakaIndia,
   author = {Mohanan, Manoj and Malani, Anup and Krishnan, Kaushik and Acharya, Anu},
   title = {Prevalence of SARS-CoV-2 in Karnataka, India},
   journal = {JAMA},
   volume = {325},
   number = {10},
   pages = {1001-1003},
   abstract = {Low- and middle-income countries contain the majority of confirmed cases of severe acute respiratory syndrome coronavirus 2 (SARS-CoV-2). India has the second highest number of reported cases, but most seroprevalence estimates come from cities. Cities, with denser population, are more vulnerable to SARS-CoV-2. However, millions of city workers fled to rural India where lockdown was less stringent.We assessed SARS-CoV-2 prevalence among volunteers from population-representative households in urban and rural areas of the state of Karnataka (population, 67.5 million).},
   ISSN = {0098-7484},
   DOI = {10.1001/jama.2021.0332},
   url = {https://doi.org/10.1001/jama.2021.0332},
   year = {2021},
   type = {Journal Article}
}

@article{Malani2021covidSeroMumbaiIndia,
   author = {Malani, Anup and Shah, Daksha and Kang, Gagandeep and Lobo, Gayatri Nair and Shastri, Jayanthi and Mohanan, Manoj and Jain, Rajesh and Agrawal, Sachee and Juneja, Sandeep and Imad, Sofia and Kolthur-Seetharam, Ullas},
   title = {Seroprevalence of SARS-CoV-2 in slums versus non-slums in Mumbai, India},
   journal = {The Lancet Global Health},
   volume = {9},
   number = {2},
   pages = {e110-e111},
   note = {doi: 10.1016/S2214-109X(20)30467-8},
   ISSN = {2214-109X},
   DOI = {10.1016/S2214-109X(20)30467-8},
   url = {https://doi.org/10.1016/S2214-109X(20)30467-8},
   year = {2021},
   type = {Journal Article}
}

@article{Imai2021twoStageExperiment,
   author = {Imai, Kosuke and Jiang, Zhichao and Malani, Anup},
   title = {Causal Inference With Interference and Noncompliance in Two-Stage Randomized Experiments},
   journal = {Journal of the American Statistical Association},
   volume = {116},
   number = {534},
   pages = {632-644},
   note = {doi: 10.1080/01621459.2020.1775612},
   ISSN = {0162-1459},
   DOI = {10.1080/01621459.2020.1775612},
   url = {https://doi.org/10.1080/01621459.2020.1775612},
   year = {2021},
   type = {Journal Article}
}


% BOOK

@book{HoldenMalani2021blockchainHoldup,
   author = {Holden, Richard and Malani, Anup},
   title = {Can Blockchain Solve the Hold-up Problem in Contracts?},
   publisher = {Cambridge University Press},
   address = {Cambridge},
   abstract = {A vexing problem in contract law is modification. Two parties sign a contract but before they fully perform, they modify the contract. Should courts enforce the modified agreement? A private remedy is for the parties to write a contract that is robust to hold-up or that makes the facts relevant to modification verifiable. Provisions accomplishing these ends are renegotiation-design and revelation mechanisms. But implementing them requires commitment power. Conventional contract technologies to ensure commitment – liquidated damages – are disfavored by courts and themselves subject to renegotiation. Smart contracts written on blockchain ledgers offer a solution. We explain the basic economics and legal relevance of these technologies, and we argue that they can implement liquidated damages without courts. We address the hurdles courts may impose to use of smart contracts on blockchain and show that sophisticated parties' ex ante commitment to them may lead courts to allow their use as pre-commitment devices.},
   ISBN = {9781009001397},
   DOI = {DOI: 10.1017/9781009004794},
   url = {https://www.cambridge.org/core/elements/can-blockchain-solve-the-holdup-problem-in-contracts/E58E833EE3090061C1C69CF05AF134E6},
   year = {2021},
   type = {Book}
}

% WORKING PAPER

@article{GuptaMalani2021inequalityIndiaCOVID,
   author = {Gupta, Arpit and Malani, Anup and Woda, Bartek},
   title = {Explaining the Income and Consumption Effects of COVID in India},
   journal = {National Bureau of Economic Research Working Paper Series},
   volume = {No. 28935},
   note = {Author contact info: Arpit Gupta E-Mail: arpit.gupta@stern.nyu.edu Anup Malani University of Chicago Law School 1111 E. 60th Street Chicago, IL 60637 Tel: 773/702-9602 E-Mail: amalani@uchicago.edu Bartosz Woda University of Chicago Law School 1111 E. 60th St. Chicago, IL 60637 USA E-Mail: bartekwoda@gmail.com},
   abstract = {The COVID-19 pandemic led to stark reductions in economic activity in India. We employ CMIE's Consumer Pyramids Household Survey to examine the timing, distribution, and mechanism of the impacts from this shock on income and consumption through December 2020. First, we estimate large and heterogeneous drops in income, with ambiguous effects on inequality. While incomes of salaried workers fell 35\%; incomes of daily laborers fell 75\%. At the same time, we observe that income fell more for individuals from households in the highest income quartile. Second, we document an increase in effort to buffer income shocks by switching occupations. We employ a Roy Model to estimate the gains from occupation churn and find, surprisingly, that reservation wages fell, implying that the risk of COVID did not reduce the value of employment. Third, we find that consumption fell less than income, suggesting households were able to smooth the idiosyncratic components of the COVID shock as well as they did before COVID. Finally, consumption of food and fuel fell less than consumption of durables such as clothing and appliances. Following Costa (2001) and Hamilton (2001), we estimate Engel curves and find that changes in consumption reflect large price shocks (rather than a retreat to subsistence) in sectors other than food and fuel/power. In the food sector, it appear that lockdown successfully distinguished essential and non-essential services, at least to the extent that it did not increase the relative price of food. There is some suggestive evidence that the price shocks outside the food sector were larger in places with greater COVID-19 cases, even during the lockdown.},
   DOI = {10.3386/w28935},
   url = {http://www.nber.org/papers/w28935},
   year = {2021},
   type = {Journal Article}
}

@article{RN6609,
   author = {Gupta, Arpit and Malani, Anup and Woda, Bartosz},
   title = {Inequality in India Declined During COVID},
   journal = {National Bureau of Economic Research Working Paper Series},
   volume = {No. 29597},
   note = {Author contact info: Arpit Gupta E-Mail: arpit.gupta@stern.nyu.edu Anup Malani University of Chicago Law School 1111 E. 60th Street Chicago, IL 60637 Tel: 773/702-9602 E-Mail: amalani@uchicago.edu Bartosz Woda University of Chicago Law School 1111 E. 60th St. Chicago, IL 60637 USA E-Mail: bartekwoda@gmail.com},
   abstract = {We use a large, representative panel data set from India with monthly data on household finances to examine the incidence of economic harms during the COVID pandemic. We observe a sharp spike in poverty, peaking during India's sharp but short lockdown. However, there was a striking decrease in income inequality outside the lockdown. There was a smaller decrease in consumption inequality, likely due to consumption smoothing. Evidence supports two mechanisms for the decline in income inequality: the capital income of top-quartile earners covaries more with aggregate income, and demand for labor fell more for higher quartiles.},
   DOI = {10.3386/w29597},
   url = {http://www.nber.org/papers/w29597},
   year = {2021},
   type = {Journal Article}
}


%%%%%%%%%%
%% 2020 %%
%%%%%%%%%%

@article{HellandMalani2020productsLiability,
   author = {Helland, Eric and Lakdawalla, Darius and Malani, Anup and Seabury, Seth A.},
   title = {Unintended Consequences of Products Liability: Evidence from the Pharmaceutical Market*},
   journal = {The Journal of Law, Economics, and Organization},
   volume = {36},
   number = {3},
   pages = {598-632},
   abstract = {We explain a surprising effect of tort liability in the market for prescription drugs. Greater punitive damage risk seems to increase prescription drug utilization in states without non-economic damage caps but decrease utilization in states with such caps. We offer an explanation for this puzzle. The vertical production process for drugs involves national upstream producers (drug companies) and local downstream producers (doctors). When a single state reallocates liability from downstream to upstream producers, national drug companies see little reason to alter their nationwide output decisions, but local physicians have incentives to increase their prescriptions in that state. The net result is higher local output. We show how this dynamic can explain our puzzle by demonstrating that punitive damages shift liability upstream from doctors to drug companies, but not when non-economic damages caps limit physician malpractice liability. We provide evidence explaining when, how, and why this type of liability shifting occurs (JEL K13, I11, I18).},
   ISSN = {8756-6222},
   DOI = {10.1093/jleo/ewaa017},
   url = {https://doi.org/10.1093/jleo/ewaa017},
   year = {2020},
   type = {Journal Article}
}

% WORKING PAPER

@article{KapoorMalani2020authoritarianGovtCovidData,
   author = {Kapoor, Mudit and Malani, Anup and Ravi, Shamika and Agrawal, Arnav},
   title = {Authoritarian governments appear to manipulate COVID data},
   journal = {arXiv preprint arXiv:2007.09566},
   year = {2020},
   type = {Journal Article}
}

@article{Malani2020adaptiveControlCOVID,
   author = {Malani, Anup and Soman, Satej and Asher, Sam and Novosad, Paul and Imbert, Clement and Tandel, Vaidehi and Agarwal, Anish and Alomar, Abdullah and Sarker, Arnab and Shah, Devavrat and Shen, Dennis and Gruber, Jonathan and Sachdeva, Stuti and Kaiser, David and Bettencourt, Luis M. A.},
   title = {Adaptive Control of COVID-19 Outbreaks in India: Local, Gradual, and Trigger-based Exit Paths from Lockdown},
   journal = {National Bureau of Economic Research Working Paper Series},
   volume = {No. 27532},
   note = {Author contact info: Anup Malani University of Chicago Law School 1111 E. 60th Street Chicago, IL 60637 Tel: 773/702-9602 E-Mail: amalani@uchicago.edu Satej Soman 1155 E 60th St Chicago, IL 60637 USA E-Mail: satej@uchicago.edu Sam Asher Johns Hopkins SAIS 1717 Massachusetts Ave NW Office 743a Washington, DC 20433 Tel: (202) 663-5830 E-Mail: samuel.asher@gmail.com Paul Novosad Economics Department Dartmouth College 6106 Rockefeller Center Room 301 Hanover, NH 03755 E-Mail: paul.novosad@dartmouth.edu Clement Imbert Department of Economics The Social Sciences Building The University of Warwick Coventry CV4 7AL United Kingdom E-Mail: C.Imbert@warwick.ac.uk Vaidehi Tandel IDFC Institute Mumbai India E-Mail: vaidehi.tandel@gmail.com Anish Agarwal 32-D666 32 Vassar Street Ray and Maria Stata Center, MIT Cambridge, MA 02139 United States of America E-Mail: anish90@mit.edu Abdullah Alomar Center for Complex Engineering Systems (CCES) King Abdulaziz City for Science and Technology (KAST) E-Mail: a.alomar@cces-kacst-mit.org Arnab Sarker Massachusetts Institute of Technology Institute for Data, Systems, and Society E-Mail: arnabs@mit.edu Devavrat Shah Massachusetts Institute of Technology E-Mail: devavrat@mit.edu Dennis Shen Massachusetts Institute of Technology E-Mail: deshen@mit.edu Jonathan Gruber Department of Economics, E52-434 MIT 50 Memorial Drive Cambridge, MA 02142 Tel: 617/253-8892 Fax: 617/253-1330 E-Mail: gruberj@mit.edu Stuti Sachdeva Massachusetts Institute of Technology E-Mail: ss5368@columbia.edu David Kaiser 182 Memorial Dr Cambridge, MA 02142 USA E-Mail: dikaiser@mit.edu Luis Bettencourt 1155 E 60th St Chicago, IL 60637 USA E-Mail: bettencourt@uchicago.edu
http://www.nber.org/papers/w27532.pdf},
   abstract = {Managing the outbreak of COVID-19 in India constitutes an unprecedented health emergency in one of the largest and most diverse nations in the world. On May 4, 2020, India started the process of releasing its population from a national lockdown during which extreme social distancing was implemented. We describe and simulate an adaptive control approach to exit this situation, while maintaining the epidemic under control. Adaptive control is a flexible counter-cyclical policy approach, whereby different areas release from lockdown in potentially different gradual ways, dependent on the local progression of the dis- ease. Because of these features, adaptive control requires the ability to decrease or increase social distancing in response to observed and projected dynamics of the disease outbreak. We show via simulation of a stochastic Susceptible-Infected-Recovered (SIR) model and of a synthetic intervention (SI) model that adaptive control performs at least as well as immediate and full release from lockdown starting May 4 and as full release from lockdown after a month (i.e., after May 31). The key insight is that adaptive response provides the option to increase or decrease socioeconomic activity depending on how it affects disease progression and this freedom allows it to do at least as well as most other policy alternatives. We also discuss the central challenge to any nuanced release policy, including adaptive control, specifically learning how specific policies translate into changes in contact rates and thus COVID-19's reproductive rate in real time.},
   DOI = {10.3386/w27532},
   url = {http://www.nber.org/papers/w27532},
   year = {2020},
   type = {Journal Article}
}



@article{Malani2020covidBiharMigrantPrevalence,
   author = {Malani, Anup and Mohanan, Manoj and Kumar, Chanchal and Kramer, Jake and Tandel, Vaidehi},
   title = {Prevalence of SARS-CoV-2 among workers returning to Bihar gives snapshot of COVID across India},
   journal = {MedRxiv},
   pages = {2020.06. 26.20138545},
   year = {2020},
   type = {Journal Article}
}

@article{HemelMalani2020immunityPassportCOVID,
   author = {Hemel, Daniel J and Malani, Anup},
   title = {Immunity passports and moral hazard},
   journal = {U of Chicago, Public Law Working Paper},
   number = {743},
   year = {2020},
   type = {Journal Article}
}


